\chapter{Emergenz der Naturkonstanten}

\section{Zwei Ebenen physikalischer Konstanten}

In der zweistufigen Theorie (IWT → WDBT+) gibt es zwei Arten von Konstanten:

\begin{enumerate}
    \item \textbf{Fundamentale IWT-Parameter:}  
    Dies sind die dimensionslosen Kopplungen und Skalen des diskreten Informationsnetzwerks. Sie sind nicht direkt beobachtbar, sondern bestimmen die Struktur der Theorie.
    
    \item \textbf{Emergente WDBT+-Konstanten:}  
    Dies sind die beobachtbaren Naturkonstanten der etablierten Physik (\( c, \hbar, G, \alpha, k_B, m_\nu, \dots \)). Sie emergieren aus den IWT-Parametern im Kontinuumslimes.
\end{enumerate}

Die Beziehung zwischen beiden Ebenen ist wechselseitig:
\begin{itemize}
    \item Aus den IWT-Parametern können die WDBT+-Konstanten berechnet werden.
    \item Umgekehrt: Aus den gemessenen WDBT+-Konstanten können die IWT-Parameter bestimmt werden.
\end{itemize}

\section{Fundamentale IWT-Parameter}

Die IWT wird durch folgende fundamentale Parameter charakterisiert:

\[
\begin{array}{ll}
\alpha_{\text{IWT}} & \text{Kopplung des lokalen (Weber-)Informationsflusses} \\
\beta_{\text{IWT}} & \text{Kopplung des globalen (Bohm-)Informationsflusses} \\
\gamma_{\text{IWT}} & \text{Kopplung der fraktalen Skalierung} \\
\lambda_0 & \text{fundamentale Längenskala des Netzwerks} \\
T & \text{fundamentaler Zeitschrift (Update-Periode)} \\
D & \text{fraktale Informationsdimension (siehe Kapitel 6)}
\end{array}
\]

Diese Parameter sind dimensionslos (bis auf \( \lambda_0 \) und \( T \), die eine absolute Skala setzen). Sie sind die „Schrauben“ der Theorie.

\section{Emergente WDBT+-Konstanten}

Im Kontinuumslimes (\( T \to 0, \lambda_0 \to 0 \), feine Diskretisierung) emergieren die bekannten Naturkonstanten:

\[
\begin{array}{ll}
c & \text{Lichtgeschwindigkeit} \\
\hbar & \text{Plancksches Wirkungsquantum} \\
G & \text{Gravitationskonstante} \\
\alpha & \text{Feinstrukturkonstante} \\
k_B & \text{Boltzmann-Konstante} \\
m_\nu & \text{Neutrinomasse (siehe Kapitel 13)}
\end{array}
\]

\section{Zusammenhang zwischen IWT-Parametern und WDBT+-Konstanten}

Aus der diskreten Informationsdynamik folgen die folgenden Beziehungen (hier ohne vollständige Herleitung, aber als Ergebnis der Emergenz):

\subsection{Lichtgeschwindigkeit}

\[
c = \frac{\lambda_0}{T}
\]

Die Lichtgeschwindigkeit ist das Verhältnis von fundamentaler Länge zu fundamentaler Zeit.

\subsection{Plancksches Wirkungsquantum}

\[
\hbar = \alpha_{\text{IWT}} \cdot \Delta I_{\text{min}} \cdot \lambda_0^2
\]

wobei \( \Delta I_{\text{min}} \) der minimale Informationsunterschied zwischen zwei Knoten ist.

\subsection{Gravitationskonstante}

\[
G = \beta_{\text{IWT}} \cdot \frac{\lambda_0^{3-D}}{T^2}
\]

Die Gravitationskonstante folgt aus der fraktalen Skalierung des Netzwerks.

\subsection{Feinstrukturkonstante}

Die Feinstrukturkonstante ist das Verhältnis lokaler zu globaler Kopplung:

\[
\alpha = \frac{\alpha_{\text{IWT}}}{\beta_{\text{IWT}}}
\]

Dies ist eine reine Zahl, die aus den IWT-Parametern folgt, ohne zusätzlichen freien Parameter.

\subsection{Boltzmann-Konstante}

\[
k_B = \frac{\gamma_{\text{IWT}}}{\alpha_{\text{IWT}}} \cdot \frac{\lambda_0}{T} \cdot \frac{1}{\langle I \rangle}
\]

wobei \( \langle I \rangle \) der mittlere Informationswert im Netzwerk ist.

\subsection{Neutrinomasse}

Die Neutrinomasse folgt aus der fundamentalen Länge (siehe Kapitel 13):

\[
m_\nu = \frac{\hbar}{\lambda_0 c}
\]

\section{Umkehrung: IWT-Parameter aus gemessenen Konstanten}

Die obigen Beziehungen können umgestellt werden, um die IWT-Parameter aus den gemessenen Naturkonstanten zu berechnen:

\[
\begin{array}{l}
\lambda_0 = \frac{\hbar}{m_\nu c} \\[4pt]
T = \frac{\lambda_0}{c} = \frac{\hbar}{m_\nu c^2} \\[4pt]
\alpha_{\text{IWT}} = \frac{\hbar}{ \Delta I_{\text{min}} \lambda_0^2 } \\[4pt]
\beta_{\text{IWT}} = \frac{G T^2}{\lambda_0^{3-D}} \\[4pt]
\gamma_{\text{IWT}} = \alpha_{\text{IWT}} \cdot \frac{k_B T}{\lambda_0} \cdot \langle I \rangle
\end{array}
\]

Die verbleibenden Größen (\( \Delta I_{\text{min}}, D, \langle I \rangle \)) sind entweder aus der Theorie bestimmbar (D) oder müssen aus zusätzlichen Bedingungen folgen.

\section{Zusammenfassung}

Die IWT und die WDBT+ sind durch eine klare Hierarchie der Konstanten verbunden:

\begin{itemize}
    \item Die IWT hat fundamentale Parameter (\( \alpha_{\text{IWT}}, \beta_{\text{IWT}}, \gamma_{\text{IWT}}, \lambda_0, T, D \)).
    \item Die WDBT+ hat emergente Konstanten (\( c, \hbar, G, \alpha, k_B, m_\nu \)).
    \item Es existieren eindeutige Beziehungen zwischen beiden Ebenen.
    \item Diese Beziehungen können umgekehrt werden, um die IWT-Parameter aus gemessenen Naturkonstanten zu bestimmen.
\end{itemize}

Damit ist die Theorie geschlossen: Keine freien Parameter in der WDBT+, alle beobachtbaren Konstanten folgen aus der IWT (oder umgekehrt). Die Feinstrukturkonstante \( \alpha \) ist kein freier Parameter mehr, sondern das Verhältnis \( \alpha_{\text{IWT}} / \beta_{\text{IWT}} \).

\section{Ausblick}

Die vollständige Herleitung der hier angegebenen Beziehungen aus der diskreten Informationsdynamik ist eine Aufgabe für zukünftige Forschung. Die vorliegende Arbeit zeigt jedoch, dass ein solcher Zusammenhang existiert und dass die Theorie konsistent formuliert werden kann.